\scanpage{170}\mspage{115}
\deleted{\unsure{correspondent à la}}\inserted{interviennent directement} quand on applique la recette générale de construction d'un dérivateur à partir d'une catégorie de modèles $(\unsure{M|\mathcal M},W)$. Donc je ne vois pas de simplification substantielle à apporter au th.-scholie 1. Tout au plus pourrait-on élaguer le diagramme des catégories envisagé, en laissant tomber $\underline{\mathrm{Liss}}\,S$, $\underline{\mathrm{Liss}}_0\,S$, $\underline{\mathrm{Prop}}\,S$, $\underline{\mathrm{Prop}}_0\,S$, qui n'y jouent pas un rôle vraiment utile.
Pour ce qui concerne les résultats énoncés, \unsure{ce|en ce} qui les concerne est couvert par ce qui est dit pour les sous-catégories pleines $\unsure{\mathcal F|F}$ de $\mathrm{Cat}_{/S}$ qui contiennent (disons) $\underline{\mathrm{Parf}}\,S$ (ou au choix, celles qui contiennent $\underline{\mathrm{Fib}}\,S$ ou $\underline{\mathrm{Cofib}}\,S$).

\figure{\begin{tikzpicture}\draw (0,0) -- (0.36,-0.66);\end{tikzpicture}}

La description des catégories du type $\mathrm{HOT}(S)$, pour $S$ fixé, me paraît à présent parfaitement bien comprise. Il est temps de faire varier $S$, et de comprendre la variation contravariante en $S$. On se donne donc
\[
S'\xrightarrow{t}S,
\]
et on se propose de décrire les foncteurs images inverses pour les trois types de catégories $\mathrm{HOT}(S)$. Mais nous allons nous borner à présent aux $\mathrm{HOT}_W(S)$ et $\mathrm{HOT}^{\mathrm{lc}}_W(S)$, puisque tout ce qui concerne $\mathrm{HOT}^{\unsure{o|0|\circ}}_W(S)=\mathrm{HOT}_W(S^o)$ \deleted{\unsure{s'en}} se déduit par dualité $S\longmapsto S^o$. Donc il s'agit de décrire
% archive sheet number: 169
